\documentclass[11pt]{article}
\usepackage[margin=1in]{geometry}
\usepackage{amsmath,amssymb}
\usepackage{booktabs}
\usepackage{graphicx}
\usepackage{array}
\usepackage{siunitx}
\usepackage[hidelinks]{hyperref}
\usepackage{caption}
\captionsetup{font=small,labelfont=bf}
\graphicspath{{outputs/}}
\newcommand{\Wj}{W_{j}}
\newcommand{\Swj}{S_{w,j}}
\sisetup{detect-weight=true,round-mode=places,round-precision=2}

\title{\textbf{A Mechanistically--Informed, Event--Based Phenomenological Model
of Wetland Effects on Flood Response} \\[4pt]
\large MMSD (Milwaukee) Watersheds --- Data, Methods, and Results}
\author{Wetland--Flood Effectiveness Study \\ \texttt{Mark's Model} pipeline}
\date{July 2026}

\begin{document} 
\maketitle

\begin{abstract}
\noindent
We operationalize the discussion note's ``mechanistically--informed phenomenological
model'' of how wetland configuration relates to flood response into an executable
event--based pipeline for nine long--record USGS gauges in the Milwaukee Metropolitan
Sewerage District (MMSD) area. Rather than annual aggregates, the unit of analysis is the
\emph{(gauge $\times$ storm event)}. For each of 2{,}699 storm events (2001--2026) we
extract hydrograph--response metrics from 15--minute streamflow, pair them with radar
(Stage~IV) event rainfall, and regress them on a cumulative wetland--effectiveness index
$\Wj=\sum_i A_i C_i M_i f(T_{ij})$, storage capacity $\Swj$, precipitation, and urbanization/
soil/terrain controls, through a progression of four nested models. Every predictor is
anchored to a published method or public dataset. With only nine (nested, confounded) gauges
we use gauge random effects, gauge--clustered standard errors, and---because nine clusters
render asymptotic standard errors unreliable---an exact wild cluster bootstrap. The wetland
\emph{level} effect is not identifiable (indistinguishable from random gauge relabeling), but
two mechanism--specific, event--identified signals survive placebo, leave--one--gauge--out,
and wild--cluster--bootstrap falsification: wetlands shave flood \emph{peaks} more strongly in
larger storms, and they \emph{widen} storm hydrographs more in larger storms. Linking discharge
to economic damages is intentionally left out of scope.
\end{abstract}

\section{Modeling framework}
Following the discussion note and the accompanying workflow, we take the
\textbf{statistical / phenomenological} route rather than a calibrated process model
(SWAT/HEC--HMS): we regress observed event--response metrics on wetland configuration and
rainfall, borrowing variable \emph{definitions} from the process literature so that no
non--authoritative parameters are invented. The full specification is
\begin{equation}
Y_{j,e} = \beta_0 + \beta_1 P_e + \beta_2 \Wj + \beta_3 (\Wj\!\cdot\!P_e)
   + \beta_4 \frac{V_e}{\Swj} + \beta_5\!\left(\Wj\!\cdot\!\tfrac{V_e}{\Swj}\right)
   + X_{j,e}\gamma + u_j + \varepsilon_{j,e},
\label{eq:full}
\end{equation}
where $Y_{j,e}$ is a response metric at gauge $j$ for storm $e$, $P_e$ storm precipitation,
$V_e$ storm rainfall volume, $X$ controls, and $u_j$ a gauge effect. We progress through
increasingly complex models to manage collinearity (Sec.~\ref{sec:models}).

\section{Study area and gauge network}
The study units are nine long--record USGS discharge gauges spanning the wetland/urbanization
gradient of the Milwaukee and Menomonee river systems (Table~\ref{tab:events}). They range
from large, slow, wetland--rich rural rivers (Milwaukee River, Cedar Creek) to small, flashy,
heavily urbanized creeks (Honey Creek, Underwood Creek). Contributing catchments, the
river network, per--cell travel--time surfaces $T_{ij}$, and conditioned DEMs are reused from
the existing routing pipeline (WhiteBox breaching + \texttt{pysheds} routing on a 10\,m,
EPSG:32616 master grid).

\section{Data and variable construction}

\subsection{Event hydrograph metrics $Y_{j,e}$ (streamflow)}
We retrieve 15--minute instantaneous discharge (USGS parameter 00060) and gage height
(00065) via \texttt{dataretrieval}, from 2001--10 (the pre--roll before the Stage~IV radar
era) to present, aggregate to hourly, and separate baseflow with the Eckhardt (2005)
recursive digital filter (hourly recession constant $a=0.98^{1/24}$, $\mathrm{BFI}_{\max}=0.60$).
Storm events are detected as prominent peaks in the quickflow series; event windows are set
at $90\%$ of peak prominence and the hydrograph width at the half--prominence full width.
Per event we compute peak discharge $Q_p$, peak stage, time--to--peak, hydrograph width,
recession constant, quickflow volume, Richards--Baker (2004) flashiness over the event
window, and hours above the gauge's 99th percentile flow. Gage height is available only
from 2017 (seven gauges) or 2007 (two gauges), so stage metrics are populated where present
and retained for the (out--of--scope) damage linkage. This yields \textbf{2{,}699 events}
(Table~\ref{tab:events}); event counts scale inversely with basin size, as expected.

\begin{table}[htbp]\centering
\caption{Step~1 --- event extraction. Median values per gauge over 2001--2026.}
\label{tab:events}
\small
\begin{tabular}{llrrrrr}
\toprule
Site & Description & $n$ events & with stage & med.\ dur.\ (h) & med.\ TTP (h) & med.\ $Q_p$ (cfs) \\
\midrule
04087000 & Milwaukee R.\ at Milwaukee   & 126 & 53  & 197 & 32.5 & 2756 \\
04086600 & Milwaukee R.\ nr Cedarburg   & 62  & 29  & 310 & 66.0 & 2145 \\
04086500 & Cedar Creek                  & 77  & 33  & 230 & 64.0 & 575  \\
04087030 & Menomonee R.\ at M.\ Falls   & 164 & 68  & 127 & 11.0 & 203  \\
04087050 & Little Menomonee nr Freistadt& 204 & 95  & 82  & 11.0 & 69   \\
04087070 & Little Menomonee at Milw.    & 322 & 273 & 59  & 8.0  & 156  \\
04087120 & Menomonee R.\ at Wauwatosa   & 504 & 196 & 35  & 3.0  & 984  \\
04087088 & Underwood Creek              & 504 & 183 & 15  & 3.0  & 281  \\
04087119 & Honey Creek                  & 736 & 535 & 10  & 2.0  & 218  \\
\bottomrule
\end{tabular}
\end{table}

\subsection{Step 1 New Model integration update}
After the baseline pipeline above was completed, I added the separate
\texttt{New Model} Step~1 event detector into the existing
\texttt{build\_events\_flow.py} script as an optional, merge-safe mode. The main change was
not to replace the original workflow, but to make the improved hydrograph detector available
inside the same Step~1 interface. Running \texttt{python build\_events\_flow.py} still uses
the original Eckhardt/quickflow method and reproduces the original Step~1 outputs. Running
\texttt{python build\_events\_flow.py --method rredi} uses the new RREDI--style detector
from \texttt{New Model/mark\_model\_pipeline.py}. This preserves the downstream contract:
Step~2 still receives one event file per gauge with the same core fields
(\texttt{site\_no}, \texttt{t\_start}, \texttt{t\_peak}, \texttt{t\_end},
\texttt{Qp\_cfs}, \texttt{quick\_vol\_m3}, \texttt{rb\_flashiness},
\texttt{peak\_month}, \texttt{peak\_year}, etc.), so the rainfall matching, panel assembly,
and model scripts do not need to be rewritten.

The input data for the new mode are the cached raw USGS discharge files created by the New
Model workflow: \texttt{usgs\_discharge\_15min\_raw/yearly/*.csv.gz}. These files contain
instantaneous discharge observations, gauge identifiers, timestamps, and source metadata.
For the handoff version I restricted the output to the same nine MMSD gauges used by the
original analysis, even though the New Model site discovery file contains a larger HUC
gauge list. I also converted the USGS local Central timestamps to UTC so that the event
windows remain compatible with the original Step~2 rainfall pairing code.

Methodologically, the new detector changes how storm events are identified. The old Step~1
method separates baseflow with the Eckhardt filter and detects prominent quickflow peaks.
The new mode instead builds a local rolling baseline directly from discharge, then defines
events by how much the observed hydrograph rises above that baseline. For each gauge, I
computed a centered 7--day rolling median baseline, formed the ratio
$R_t=Q_t/\widetilde{Q}_t$, started an event when $R_t\ge1.20$, ended the event only after
$R_t\le1.05$ for six continuous hours, merged events separated by less than 12 hours, and
kept only events whose peak ratio was at least 1.30. This produces more candidate events
than the original quickflow-prominence detector, especially in small flashy urban basins,
but it keeps the same downstream event-level structure.

\begin{table}[htbp]\centering
\caption{Variables added or reconstructed in the Step~1 New Model integration. These
variables are produced before rainfall is attached; rainfall variables such as $P_e$ and
$V_e$ are still computed later in Step~2.}
\label{tab:rredi-vars}
\small
\begin{tabular}{p{0.27\textwidth}p{0.66\textwidth}}
\toprule
Variable & How it was constructed in the integrated RREDI mode \\
\midrule
\texttt{baseline\_at\_peak\_cfs} & Centered 7--day rolling median discharge
$\widetilde{Q}_t$ evaluated at the event peak. This is the local hydrograph baseline used
instead of the Eckhardt baseflow estimate. \\
\texttt{rredi\_peak\_ratio} & Event peak discharge divided by the rolling baseline at the
peak, $Q_p/\widetilde{Q}_{p}$. Events are retained only when this ratio is at least 1.30. \\
\texttt{t\_start}, \texttt{t\_end} & Start is the last low-flow time before the hydrograph
crosses $Q/\widetilde{Q}\ge1.20$. End is the first time after the peak when
$Q/\widetilde{Q}\le1.05$ and remains below that threshold for six hours. \\
\texttt{t\_peak}, \texttt{Qp\_cfs} & The timestamp and value of the maximum observed
discharge within the event window. \\
\texttt{Qp\_quick\_cfs} & Peak excess discharge above the rolling baseline,
$\max(Q_t-\widetilde{Q}_t,0)$, within the event window. \\
\texttt{dur\_hr}, \texttt{time\_to\_peak\_hr}, \texttt{recession\_time\_hr} & Event duration,
start-to-peak time, and peak-to-end time, computed directly from the event timestamps. \\
\texttt{hydro\_width\_hr} & Width of the event hydrograph above half of the peak excess
discharge. This mirrors the original hydrograph-width concept but uses the rolling-baseline
excess series. \\
\texttt{quick\_vol\_m3} & Trapezoidal integral of excess discharge above the rolling
baseline over the event window, converted from cubic-feet-seconds to cubic meters using
$1~\mathrm{ft}^3=0.0283168~\mathrm{m}^3$. \\
\texttt{total\_vol\_m3} & Trapezoidal integral of total observed discharge over the event
window, converted to cubic meters using the same unit conversion. \\
\texttt{rb\_flashiness} & Richards--Baker style event flashiness, computed as
$\sum|\Delta Q|/\sum Q$ over finite discharge observations in the event window. \\
\texttt{recession\_k\_per\_hr} & Slope of $\log(Q)$ on the falling limb after the peak,
reported as a positive recession rate when enough valid points are available. \\
\texttt{dur\_above\_q99\_hr} & Number of event-window observation intervals above the
gauge-specific 99th percentile discharge threshold, converted to hours. \\
\texttt{rredi\_data\_coverage}, \texttt{rredi\_max\_gap\_hours},
\texttt{rredi\_qc\_pass} & Quality-control diagnostics. Coverage compares observed points
with the expected count from the median sampling interval; maximum gap flags large
streamflow gaps; an event passes QC when coverage is at least 80\%, the maximum gap is no
more than two hours, and the event has a positive recession limb. \\
\bottomrule
\end{tabular}
\end{table}

Because the New Model discharge cache contains discharge but not gage height,
\texttt{peak\_stage\_ft} is intentionally left missing in the new mode. I also added a
parquet-first output writer: if \texttt{pyarrow} or \texttt{fastparquet} is available, the
script writes the exact downstream file \texttt{data/events\_\{site\}.parquet}; otherwise it
writes a compressed CSV fallback for review. In the handoff package, I included the modified
script, README, manifest, raw discharge cache, New Model reference outputs, and validation
outputs under \texttt{step1\_new\_model\_handoff/}.

\begin{table}[htbp]\centering
\caption{Step~1 integration validation --- RREDI--style event extraction, 2001--2023.
The regression results elsewhere in this report still correspond to the original
2{,}699--event panel unless Steps~2--6 are rerun using these updated Step~1 event files.}
\label{tab:rredi-update}
\small
\begin{tabular}{lrrrrr}
\toprule
Site & raw discharge rows & events & QC--pass & med.\ dur.\ (h) & med.\ $Q_p$ (cfs) \\
\midrule
04087000 & 606{,}667  & 509  & 478  & 31.8 & 786.0 \\
04086500 & 727{,}894  & 384  & 353  & 66.8 & 138.0 \\
04086600 & 1{,}013{,}979 & 281 & 274 & 55.8 & 386.0 \\
04087120 & 855{,}658  & 1{,}059 & 954  & 47.1 & 398.0 \\
04087030 & 1{,}198{,}768 & 798 & 737 & 56.0 & 60.2 \\
04087050 & 895{,}212  & 543  & 506  & 52.2 & 11.5 \\
04087070 & 595{,}328  & 1{,}019 & 987  & 46.0 & 39.0 \\
04087088 & 1{,}984{,}220 & 1{,}368 & 1{,}303 & 43.6 & 61.8 \\
04087119 & 1{,}255{,}953 & 1{,}394 & 1{,}266 & 39.6 & 46.0 \\
\midrule
\textbf{Total} & 9{,}133{,}679 & 7{,}355 & 6{,}858 & --- & --- \\
\bottomrule
\end{tabular}
\end{table}

For handoff, the integrated script, README, manifest, raw discharge cache, New Model
reference outputs, and validation outputs were packaged in
\texttt{step1\_new\_model\_handoff/}. The package is designed so a collaborator can review
the change independently, install a parquet engine such as \texttt{pyarrow}, run the new
Step~1 mode, and then continue with the existing Step~2--Step~6 scripts.

\subsection{Event rainfall $P_e$, $V_e$, API (radar QPE)}
Basin--average storm rainfall is derived from NCEP Stage~IV radar--gauge precipitation
(1997+, $0.125^{\circ}\approx12$\,km) served as daily analyses by the Iowa Environmental
Mesonet (IEMRE) web service. For each catchment we place uniform sample points, snap them to
the analysis lattice, and area--weight the nodes; nodes are cached and shared across
catchments. Per event we obtain the window rainfall depth $P_e$, volume $V_e=P_e\cdot\text{area}$,
the event runoff coefficient (quickflow volume $/\,V_e$), storm/peak--stage response ratios,
a 30--day antecedent precipitation index ($k=0.9$), and the peak daily intensity. The
basin--mean climatology (887\,mm\,yr$^{-1}$ at the main gauge) and event runoff coefficients
(median 0.16--0.32) validate the extraction (Table~\ref{tab:rain}). The 12\,km grid smears
convective cells over the smallest ($<50$\,km$^2$) basins, where a single node covers the
catchment---a documented limitation, upgradeable to native 4\,km Stage~IV or 1\,km MRMS.

\begin{table}[htbp]\centering
\caption{Step~2 --- radar event rainfall. RC = event runoff coefficient.}
\label{tab:rain}
\small
\begin{tabular}{lrrrrrr}
\toprule
Site & area (km$^2$) & nodes & $n$ events & med.\ $P_e$ (mm) & med.\ RC & med.\ API (mm) \\
\midrule
04087000 & 1751 & 25 & 126 & 48.2 & 0.218 & 19.4 \\
04086500 & 321  & 6  & 77  & 71.0 & 0.317 & 18.6 \\
04086600 & 1518 & 21 & 62  & 75.4 & 0.267 & 17.1 \\
04087120 & 310  & 6  & 504 & 21.8 & 0.181 & 20.7 \\
04087030 & 85   & 3  & 164 & 46.0 & 0.266 & 20.7 \\
04087050 & 21   & 1  & 204 & 29.5 & 0.314 & 21.0 \\
04087070 & 50   & 2  & 322 & 27.7 & 0.289 & 20.8 \\
04087088 & 47   & 3  & 504 & 20.4 & 0.161 & 20.7 \\
04087119 & 27   & 1  & 736 & 16.0 & 0.254 & 19.3 \\
\bottomrule
\end{tabular}
\end{table}

\subsection{Cumulative wetland effectiveness $\Wj$ and storage $\Swj$}
Wetland polygons are the USFWS National Wetlands Inventory (NWI), retrieved by tiled bounding
--box queries from the FWS ArcGIS service (tiling avoids server timeouts on large basins),
deduplicated and clipped to each catchment (41{,}092 polygons). For each wetland $i$ we compute
area $A_i$; travel time $T_{ij}$ to the gauge (sampled from the routing surface $T_{j}$);
distance to the channel network ($\mathrm{acc}>0.5$\,km$^2$), giving a near--channel
connectivity indicator $C_i$ ($<100$\,m); and a Cowardin type modifier $M_i$ (emergent marsh
$1.0$; forested/scrub--shrub $0.8$; open water/riverine $0.5$; lacustrine $0.3$). The
cumulative effectiveness is
\begin{equation}
\Wj=\sum_{i} A_i\, C_i\, M_i\, f(T_{ij}), \qquad
f(T)=e^{-T/\tau}\ \text{or}\ \frac{1}{1+T/\tau},
\end{equation}
with $\tau$ set to the catchment--median travel time (absolute $T$ is uncalibrated, so a fixed
$\tau$ would be degenerate; $\Wj$ is therefore a \emph{relative}, cross--basin--comparable
index). Ablation variants (dropping $C$, dropping $M$, harmonic decay, graded $C$) are retained.
Storage uses the published Volume--Area power law $V_i=cA_i^{\beta}$ ($c=0.05$, $\beta=1.2$;
order--of--magnitude, pending a Wisconsin--specific fit); a DEM depression--volume cross--check
is uninformative here because the conditioned DEM was purpose--built to remove depressions.
The NWI wetland area at the main gauge (324\,km$^2$) matches the independent NLCD estimate
(325\,km$^2$). The wetland gradient is exactly the one the analysis needs: rural basins
18--22\% wetland, urban creeks 1--11\% (Table~\ref{tab:wsw}).

\begin{table}[htbp]\centering
\caption{Step~3 --- NWI wetland effectiveness and storage. $W_{\exp}$ in m$^2$; $S_w$ Volume--Area.}
\label{tab:wsw}
\small
\begin{tabular}{lrrrrrr}
\toprule
Site & $n$ wet & wet (km$^2$) & wet frac & near frac & $W_{\exp}$ ($10^6$) & $S_w$ depth (mm) \\
\midrule
04087000 & 16826 & 323.7 & 0.185 & 0.52 & 51.4 & 94.9 \\
04086500 & 3710  & 71.2  & 0.222 & 0.48 & 12.0 & 123.7 \\
04086600 & 14666 & 304.6 & 0.201 & 0.53 & 51.1 & 103.2 \\
04087120 & 2802  & 33.1  & 0.107 & 0.54 & 5.3  & 50.1 \\
04087030 & 1016  & 16.2  & 0.190 & 0.44 & 2.5  & 95.0 \\
04087050 & 176   & 2.4   & 0.112 & 0.53 & 0.5  & 53.9 \\
04087070 & 439   & 5.3   & 0.106 & 0.62 & 1.3  & 48.8 \\
04087088 & 375   & 4.0   & 0.084 & 0.57 & 0.7  & 38.7 \\
04087119 & 82    & 0.3   & 0.012 & 0.71 & 0.06 & 3.7 \\
\bottomrule
\end{tabular}
\end{table}

\subsection{Control variables}
To separate the wetland signal from urbanization, soils, terrain, and size we compute per
catchment: drainage area $\mathrm{DA}$; basin and channel slope from the 3DEP DEM; NLCD
impervious fraction (2001/2011/2021, interpolated to each event's year); POLARIS surface
$K_{\text{sat}}$ and sand fraction; and gNATSGO available water storage. Impervious cover
spans 4.3\% (rural Cedar Creek) to 52.9\% (urban Honey Creek) and trends upward through time;
urban basins also have lower $K_{\text{sat}}$, less sand, and steeper channels
(Table~\ref{tab:ctrl}).

\begin{table}[htbp]\centering
\caption{Step~4 --- control variables. $K_{\text{sat}}$ in $\log_{10}$(cm\,h$^{-1}$).}
\label{tab:ctrl}
\small
\begin{tabular}{lrrrrrrr}
\toprule
Site & DA (km$^2$) & basin slope & chan slope & imp'01 (\%) & imp'21 (\%) & $K_{\text{sat}}$ & AWS (mm) \\
\midrule
04087000 & 1751 & 0.047 & 0.013 & 6.7  & 7.6  & 0.380 & 164 \\
04086500 & 321  & 0.040 & 0.012 & 3.4  & 4.8  & 0.372 & 176 \\
04086600 & 1518 & 0.049 & 0.012 & 3.5  & 4.3  & 0.383 & 167 \\
04087120 & 310  & 0.038 & 0.020 & 26.8 & 29.7 & 0.333 & 159 \\
04087030 & 85   & 0.034 & 0.013 & 13.6 & 16.4 & 0.337 & 180 \\
04087050 & 21   & 0.033 & 0.013 & 4.2  & 5.6  & 0.316 & 169 \\
04087070 & 50   & 0.034 & 0.016 & 19.0 & 20.8 & 0.328 & 154 \\
04087088 & 47   & 0.043 & 0.024 & 31.0 & 33.5 & 0.351 & 154 \\
04087119 & 27   & 0.033 & 0.030 & 51.6 & 52.9 & 0.296 & 144 \\
\bottomrule
\end{tabular}
\end{table}

\subsection{Panel assembly and collinearity}
The four blocks are joined into a single $(\text{gauge}\times\text{event})$ table of
\textbf{2{,}699 rows (2{,}685 usable}, with real storm rainfall and finite runoff
coefficient). The gauge--level predictors are strongly collinear
(Fig.~\ref{fig:collin}): wetland fraction versus impervious cover $r=-0.87$; the extensive
$W_{\exp}$ is nearly identical to storage $S_w$ ($r=1.0$) and highly correlated with drainage
area ($r=0.89$); impervious versus channel slope $r=0.97$. Two consequences drive the
modeling: (i) we use the \emph{intensive} treatment $W_n=W_{\exp}/\text{area}$ (an
effectiveness--weighted wetland fraction, less area--collinear) and control $\mathrm{DA}$
explicitly; (ii) because $W_n$ is gauge--constant, the wetland \emph{level} effect is a
purely cross--sectional nine--point contrast confounded with urbanization, whereas the
event--varying interaction and saturation terms carry within--gauge identification.

\begin{figure}[htbp]\centering
\includegraphics[width=0.62\textwidth]{step5_collinearity.png}
\caption{Gauge--level predictor correlations ($n=9$). Wetland measures are strongly
anti--correlated with impervious cover and channel slope, and $W_{\exp}$ tracks storage and
drainage area.}
\label{fig:collin}
\end{figure}

\section{Statistical models}\label{sec:models}
All predictors are standardized. We fit the nested progression of Eq.~\eqref{eq:full}:
\begin{itemize}\itemsep2pt
\item \textbf{M1 (basic wetland effect):} $Y\sim P + W + X$. \emph{Do more effective wetlands
reduce peaks?} ($\beta_2$)
\item \textbf{M2 (peak--shaving):} $+\,W\!\cdot\!P$. \emph{Do wetlands matter more in large
storms?} ($\beta_3<0$)
\item \textbf{M3 (saturation):} $+\,\log(V_e/S_w)$. \emph{Does response rise once storm
volume exceeds storage?}
\item \textbf{M4 (nonlinear saturation):} $+\,W\!\cdot\!(V_e/S_w)$. \emph{Does wetland
effectiveness decline past capacity?}
\end{itemize}
Controls $X$ are $\log\mathrm{DA}$, impervious, $K_{\text{sat}}$, channel slope, API, and a
winter flag. Each model is estimated with (a) a gauge random--intercept mixed model and
(b) OLS with gauge--clustered standard errors. Outcomes are log peak discharge (primary),
Richards--Baker flashiness, and log hydrograph width. Falsification follows the workflow:
variance inflation factors, a placebo that permutes wetlands across gauges, and
leave--one--gauge--out refitting.

\section{Results}

\subsection{Coefficients and asymptotic diagnostics}
For log peak discharge, precipitation dominates ($\hat\beta_P\approx0.42$ per SD, a sanity
check). The wetland \emph{level} effect $\hat\beta_2$ is essentially zero and insignificant
in every model; the \textbf{peak--shaving interaction $W\!\cdot\!P$ is negative and stable}
(Table~\ref{tab:coef}, Fig.~\ref{fig:forest}). Variance inflation confirms the identification
structure: the between--gauge terms are severely collinear ($\mathrm{VIF}$: channel slope 175,
impervious 76, saturation 72, $W$ 50), whereas the within--gauge interaction $W\!\cdot\!P$ has
$\mathrm{VIF}=5$. The placebo (permuting wetlands among gauges, 500 draws) cannot distinguish
the observed $W$ effect from random relabeling ($p=0.85$) but flags the interaction as special
($p=0.006$); leave--one--gauge--out gives an unstable, sign--flipping $W$ (range
$[-0.69,+0.22]$) but a rock--stable $W\!\cdot\!P$ (range $[-0.174,-0.142]$).

\begin{table}[htbp]\centering
\caption{Step~6 --- MixedLM (gauge random effect) standardized coefficients, primary outcome
$\log Q_p$.}
\label{tab:coef}
\small
\begin{tabular}{lrrrrr}
\toprule
Model & $P$ & $W$ & $W\!\cdot\!P$ & $V_e/S_w$ & $W\!\cdot\!(V_e/S_w)$ \\
\midrule
M1 basic        & 0.421 & \phantom{-}0.044\,(0.89) & --- & --- & --- \\
M2 peak--shave  & 0.413 & $-$0.020\,(0.77) & $-$0.153\,(0.00) & --- & --- \\
M3 saturation   & 0.413 & $-$0.026\,(0.70) & $-$0.153\,(0.00) & 0.000\,(1.00) & --- \\
M4 nonlin.\ sat.& 0.525 & $-$0.072\,(0.34) & $-$0.113\,(0.00) & $-$0.164\,(0.07) & $-$0.068\,(0.04) \\
\bottomrule
\end{tabular}
\\[3pt]
\footnotesize $p$--values in parentheses (asymptotic; see wild bootstrap below for honest
few--cluster inference).
\end{table}

\begin{figure}[htbp]\centering
\includegraphics[width=\textwidth]{step6_forest.png}
\caption{Models~1--4 key coefficients (MixedLM); red = asymptotic $p<0.05$. The peak--shaving
interaction $W\!\cdot\!P$ is consistently negative for peak discharge and positive for
hydrograph width.}
\label{fig:forest}
\end{figure}

\subsection{Wild cluster bootstrap (honest few--cluster inference)}
With only nine clusters, cluster--robust asymptotic standard errors are unreliable. We apply
the Cameron--Gelbach--Miller wild cluster bootstrap, enumerating \emph{all} $2^9=512$
Rademacher sign vectors exactly (restricted bootstrap--$t$ for $p$--values; unrestricted for
confidence intervals). This is decisive: it prunes false positives that the asymptotic
$t$--statistics had inflated (e.g.\ the flashiness $W\!\cdot\!P$ moves from $p=0.005$ to
$p=0.90$; the primary $W\!\cdot\!P$ $t=-17$ corresponds to an honest $p=0.004$). Two
mechanism--specific findings survive (Table~\ref{tab:wcb}):
\begin{enumerate}\itemsep2pt
\item \textbf{Peak discharge: $W\!\cdot\!P$, $p=0.004$}, 95\% CI $[-0.169,-0.137]$ ---
wetlands attenuate flood peaks more strongly in larger storms.
\item \textbf{Hydrograph width: $W\!\cdot\!P$ ($p=0.025$) and saturation $V_e/S_w$ ($p=0.031$)}
--- wetlands widen storm hydrographs more in larger storms, and hydrograph shape responds to
the storm--to--storage ratio.
\end{enumerate}
All flashiness effects, all Model~4 nonlinear--saturation terms, and every wetland \emph{level}
effect fail to survive the bootstrap.

\begin{table}[htbp]\centering
\caption{Wild cluster bootstrap (512 exact Rademacher sign vectors, $G=9$). $p_{\text{WCR}}$
is the restricted bootstrap--$t$ $p$--value; CI is the unrestricted percentile interval.}
\label{tab:wcb}
\small
\begin{tabular}{llrrrl}
\toprule
Outcome & Term & $\hat\beta$ & cluster $t$ & $p_{\text{WCR}}$ & 95\% CI \\
\midrule
$\log Q_p$      & $W\!\cdot\!P$ (M2)          & $-$0.153 & $-$17.1 & \textbf{0.004} & $[-0.169,-0.137]$ \\
$\log Q_p$      & $W$ (M1)                    & $-$0.209 & $-$1.5  & 0.406 & $[-0.439,\,0.022]$ \\
$\log Q_p$      & $W\!\cdot\!(V_e/S_w)$ (M4)  & $-$0.067 & $-$2.7  & 0.133 & $[-0.112,-0.022]$ \\
hydro.\ width   & $W\!\cdot\!P$ (M2)          & $\phantom{-}$0.238 & 4.8 & \textbf{0.025} & $[0.149,\,0.328]$ \\
hydro.\ width   & $V_e/S_w$ (M3)              & $\phantom{-}$0.759 & 5.2 & \textbf{0.031} & $[0.508,\,1.010]$ \\
hydro.\ width   & $W$ (M1)                    & $\phantom{-}$1.794 & 4.4 & 0.100 & $[1.052,\,2.535]$ \\
flashiness      & $W\!\cdot\!P$ (M2)          & $\phantom{-}$0.002 & 0.2 & 0.897 & $[-0.013,\,0.017]$ \\
flashiness      & $V_e/S_w$ (M3)              & $\phantom{-}$0.068 & 4.5 & 0.205 & $[0.042,\,0.093]$ \\
\bottomrule
\end{tabular}
\end{table}

\section{Discussion}
The results align closely with the discussion note's hypotheses while remaining honest about
what nine confounded gauges can support. The naive question ``do wetland--rich basins have
lower peaks?'' cannot be answered here: the wetland level effect is collinear with
urbanization ($r=-0.87$), and every falsification test---placebo, leave--one--gauge--out, and
the wild cluster bootstrap---marks it as unidentified. The \emph{event--identified} mechanism
tests, however, are robust and physically coherent. The negative $W\!\cdot\!P$ on peak
discharge is precisely the peak--shaving hypothesis ($\beta_3<0$): the marginal peak per unit
rainfall is smaller in wetland--rich basins, and disproportionately so as storms grow. The
positive $W\!\cdot\!P$ and saturation effects on hydrograph width express the same physics in
the time domain---wetlands store and release water more slowly, broadening the hydrograph---
consistent with the note's expectation that wetlands affect \emph{shape and timing} more than
runoff \emph{volume} (indeed no robust runoff--coefficient or flashiness effect emerges).

\section{Limitations}
(i) Nine nested, strongly confounded gauges: only within--gauge, event--varying terms are
identified; the wetland level effect is not. (ii) Nine clusters: even the bootstrap is weak
for cluster--invariant regressors; inference is trustworthy for the interaction/saturation
terms. (iii) Rainfall at $12$\,km smears convective storms over the smallest basins.
(iv) Storage $S_w$ uses literature--range Volume--Area constants (order of magnitude) and
absolute travel times are uncalibrated, so $\Wj$ is a relative index. (v) The analysis is
associational, not causal. (vi) Engineered storage (detention basins, the Deep Tunnel) and
dam removals are not yet explicitly modeled.

\section{Conclusions and next steps}
An end--to--end event--based pipeline (nine gauges, 2{,}699 storms, 2001--2026) reproduces the
mechanistically--informed model and delivers two falsification--surviving findings: wetlands
shave flood peaks and widen hydrographs more strongly in larger storms. To move from a
robust \emph{method demonstration} toward an identifiable wetland \emph{level} effect, the
decisive next step is to expand to a regional sample of many \emph{independent} gauged
catchments spanning a wetland gradient not collinear with urbanization; secondary refinements
include native 4\,km/1\,km radar for the small basins, NOAA Atlas~14 return periods, a
Wisconsin--specific storage calibration, and a wild cluster bootstrap on the enlarged sample.
Linking discharge/stage to flood damages (depth--damage or NFIP claims) is deliberately left
to the study team.

\paragraph{Reproducibility.} The pipeline is six scripts run in sequence:
\texttt{build\_events\_flow.py} (events), \texttt{build\_events\_rain.py} (radar rainfall),
\texttt{build\_wetland\_nwi.py} ($W$, $S_w$), \texttt{build\_controls.py} (controls),
\texttt{build\_panel.py} (panel), \texttt{fit\_models.py} (Models 1--4), plus
\texttt{wild\_bootstrap.py}. All data sources are public: USGS NWIS, IEM Stage~IV, USFWS NWI,
NLCD, POLARIS, gNATSGO, and USGS 3DEP.

\end{document}
